% ==============================================================
%  Chapter 1 -- Introduction and State of the Art
% ==============================================================

\chapter{Introduction}
\label{ch:introduction}

This introductory chapter establishes the context, motivation, and core objectives of the thesis. It presents Influence Diagrams as the foundational framework for decision support systems and identifies the exhaustive parameter elicitation process as their primary practical bottleneck. To address this limitation, the chapter outlines the main proposal of this research: automating the recovery of quantitative parameters from a reduced set of expert decision rules using evolutionary optimization. Finally, it defines the scope of the study and provides a structural outline of the subsequent chapters.

% --------------------------------------------------------------
\section{Problem Definition}
% --------------------------------------------------------------

A \emph{Decision Support System} (DSS) is an interactive
computational tool designed to assist a decision-maker in the analysis
of unstructured problems by combining data, formal models and
analytical techniques~\cite{insua2005}. Within
this broad family, \emph{Influence Diagrams}~(IDs)
\cite{howard2005influence,koller2009pgm,shenoy2009eolss} have become
the de-facto standard for representing decision
problems under uncertainty and offer a particularly compact
graphical formalism:
a single directed acyclic graph encodes simultaneously the
probabilistic structure of the problem (chance nodes), the
controllable choices of the agent (decision nodes) and their
preferences over outcomes (utility nodes). Unlike a fully expanded
decision tree, an ID allows the analyst to reason about the problem
at the level of variables and to compute the optimal policy by direct
evaluation of the graph using exact algorithms such as the one of
Shachter~\cite{shachter1986evaluating}.

The expressive power of an ID rests on two distinct ingredients.
The first is its \emph{qualitative structure}, the set of nodes and
arcs, which captures the conditional dependencies of the problem.
The second is its \emph{quantitative parameterisation}, consisting of
the \emph{conditional probability tables} (CPTs) of all chance nodes and the
numerical entries of the \emph{utility function}. While the structure can
often be elicited from a domain expert in a relatively small number of
modelling sessions, the parameterisation grows exponentially with the
connectivity of the graph and becomes the real bottleneck of any
realistic ID-based DSS~\cite{bielza2010modeling}. A well-known
illustration is the neonatal-jaundice ID of G{\'o}mez et
al.~\cite{gomez2007jaundice}: more than 8\,000 probability entries
and over a hundred utility values had to be elicited from clinicians
before the model could be deployed, and the very same effort is
documented as the dominant cost of the project in the companion
analysis of Bielza et al.~\cite{bielza2000issues}.

To put these numbers in perspective, the SRI (Stanford Research Institute) elicitation
protocol, still the reference in the field, requires roughly \emph{thirty minutes per parameter}~\cite{bielza2010modeling},
which makes the direct estimation of a model with thousands of
entries impractical. The community has responded with a battery of
techniques to reduce the parameter count, divorcing parents,
sequential refinements, canonical Noisy-OR / Noisy-AND
models, and the same authors report that on the jaundice ID itself
Noisy-OR achieves a \emph{77\,\% global reduction} of the
probabilistic parameters and up to 97\,\% on individual
distributions~\cite{bielza2010modeling}. The utility node is in
general even harder to elicit: a single utility function like the one used in the neonatal-jaundice ID~\cite{gomez2007jaundice} with five
attributes already requires of the order of 5\,400 entries and the
corresponding pairwise-comparison protocol would demand more than
14 million expert judgements~\cite{bielza2010modeling}. These
figures, far from being anecdotal, are what makes the
parameter-elicitation problem the dominant cost of any ID-based
DSS deployment.

This thesis addresses precisely this parameter-elicitation problem.
We assume that the structure of the ID is given, that all variables
are discrete and that the model contains a single value node. Under
these assumptions, the question we tackle can be stated as follows:

\begin{quote}\itshape
Given the qualitative structure of an Influence Diagram and a small
set of decision rules provided by a domain expert, can we
automatically recover a parameter vector (CPTs and utilities) that
reproduces those rules, thereby relieving the expert from the burden
of specifying every numerical entry of the model?
\end{quote}


\section{Motivation}
\label{sec:motivation}

Several factors make this problem both timely and important.

\begin{itemize}
  \item \textbf{The elicitation bottleneck.} Direct elicitation techniques
        (reference lotteries, indifference judgements, pairwise
        comparisons\ldots) remain the de-facto standard for the construction of
        probabilistic decision models~\cite{bielza2010modeling}. They demand
        long interviews with scarce experts, are cognitively demanding for the
        interviewee and are well known to introduce inconsistencies that require
        subsequent revision. In domains such as clinical medicine, public-health
        planning or safety-critical engineering, the cost of expert time is the
        dominant factor preventing the wider adoption of ID-based decision
        tools.

  \item \textbf{Scarcity of historical data.} A natural alternative to direct
        elicitation is to learn the parameters of the model from historical data
        via maximum-likelihood or expectation--maximisation estimators, as
        originally proposed by Ezawa and Gupta in the context of
        IDs~\cite{ezawa1997learning}. However, the regime in which most clinical
        and engineering problems live is precisely the opposite of what such
        statistical methods require: there are very few annotated historical
        cases, but a domain expert is usually able to verbalise a small number
        of high-level rules of the form ``\emph{if condition $\mathbf{x}$ holds,
        then take action $a$}''.

  \item \textbf{Decisions instead of probabilities.} Eliciting rules is
        cognitively much closer to clinical or managerial practice than
        eliciting numerical probabilities. Practitioners reason naturally in
        terms of recommended actions under prototypical scenarios, and reach
        agreement on rules far more easily than on probability tables.
        Exploiting rules as a supervisory signal therefore aligns the
        elicitation process with the way experts actually think.


\item \textbf{Use cases.} Concrete domains where the proposed approach
is of immediate practical interest include:
\begin{itemize}
    \item \textbf{Clinical decision support.}
          Risk stratification, triage, treatment selection or the
          management of chronic diseases. The neonatal-jaundice ID
          of~\cite{gomez2007jaundice} is a canonical example, and
          the demand for evidence-based decision aids in clinical
          practice is by now well documented~\cite{dowding2004evidence}.
    \item \textbf{Industrial maintenance and safety.}
          Inspection scheduling, replacement-vs-repair choices, and
          safety-critical operations in which expert rules are
          available but historical incident data are
          sparse~\cite{machado2016maintenance}.
    \item \textbf{Public policy and emergency response.}
          Vaccination campaigns, evacuation plans and resource
          allocation under uncertainty, where domain rules exist but
          full probabilistic models do
          not~\cite{kailiponi2010evacuation}.
    \item \textbf{Cyber-security and fraud detection.}
          Operational analysts maintain rule-based playbooks that can
          be turned into the supervisory signal of an ID-based
          decision module~\cite{chockalingam2023cyber}.
\end{itemize}\end{itemize}

\section{Scope}

The field of decision support systems is vast. To keep the
contribution focused, the work in this thesis is restricted to:
\begin{enumerate}
    \item Influence Diagrams with \emph{discrete} chance and decision
          variables;
    \item a \emph{single} value (utility) node on a normalized scale of utilities $[0,10]$;
    \item a \emph{pre-specified} dependency graph (the structure is
          provided as an input and is not subject to learning);
    \item supervision signals consisting of a finite set of
          \emph{partial decision rules} produced by a domain expert.
    \item \emph{non-parametric} characteristics of the probability distributions; 
          for example, $0 < p_i < 1$ for all $i$, skewness and monotonicity, in 
          ordinal variables (low, medium, high) such as $p_{\text{low}} < p_{\text{med}} < p_{\text{high}}$ 
          or $p_{\text{low}} > p_{\text{med}} > p_{\text{high}}$, mode identification, 
          and other bounds like $a < p_i < b$.
\end{enumerate}

This input to the EDA is, in fact, a future line of research, and it should partially resolve the ill-conditioning of the problem that yields deficient or degenerate solutions.

A wealth of representational extensions has been proposed in the ID literature: asymmetric decision problems (conditional
distribution trees, sequential decision diagrams, sequential
valuation networks), continuous and hybrid IDs (Gaussian IDs,
mixtures of truncated exponentials and polynomials), unconstrained
IDs and multi-agent extensions (MAIDs). The review of Bielza,
G{\'o}mez and Shenoy~\cite{bielza2011review} offers a thorough
catalogue of them. Although this thesis deliberately restricts
itself to the regular, discrete, single-value setting described
above, the extensions are revisited only in the perspectives
chapter.

\section{Deep Learning and Neural Networks}

It is impossible to discuss modern artificial intelligence without acknowledging the profound impact of \emph{Deep Learning (DL)} and \emph{Neural Networks}. These architectures currently dominate the landscape of unstructured data processing, achieving unprecedented success in computer vision, natural language processing, and pattern recognition. However, despite their overwhelming popularity, they present severe limitations when applied to the types of decision-support scenarios targeted in this thesis.

\begin{itemize}
    \item First and most important is the issue of \emph{auditability and interpretability}. Neural networks function inherently as ``black boxes''; their knowledge is distributed across millions or billions of meaningless weights. In high-stakes domains such as clinical triage or safety-critical maintenance, a decision-maker cannot blindly trust an automated output. Influence Diagrams, in contrast, are inherently transparent. Every node, arc, and utility value possesses a clear, auditable semantic meaning that experts can inspect, debate, and mathematically validate.

    \item Secondly, Deep Learning exhibits an extreme \emph{hunger for data}. Training a robust neural network requires vast amounts of labeled historical records to avoid overfitting. As discussed in Section~\ref{sec:motivation}, the operational reality of many critical domains is characterized by a scarcity of historical data but an availability of expert heuristics. Neural networks are not suited for this data-deprived environment, whereas our proposed approach is explicitly designed to exploit these scarce but high-value expert rules.

    \item Finally, standard deep models predominantly solve prediction or classification tasks; they do not naturally embed a decision-theoretic framework. An ID explicitly separates the probabilistic beliefs (chance nodes) from the agent's interventions (decision nodes) and their preferences (utility nodes), providing a rigorous mechanism for rational decision-making under uncertainty. 
\end{itemize}

For these reasons, this thesis deliberately steers away from the Deep Learning paradigm. We choose to retain the transparency, rigorous uncertainty management, and use the Influence Diagrams to solve their only major practical drawback: the parameter-elicitation bottleneck.


% --------------------------------------------------------------
\section{Objectives and Contributions}
% --------------------------------------------------------------

The overall objective of this thesis is to develop and validate a
proof-of-concept system that estimates the quantitative parameters of
a regular Influence Diagram from a partial expert rule base, using
Estimation of Distribution Algorithms as the search engine. To
achieve this global goal we set the following specific objectives:
\begin{enumerate}
    \item Formalise the parameter-recovery problem as a constrained
          black-box optimisation over the joint space of CPTs and
          utilities.
    \item Design an encoding of the parameter vector compatible with
          continuous EDAs and capable of guaranteeing the validity of
          every decoded CPT.
    \item Implement the resulting system on top of the
          \textsc{EDAspy} library~\cite{soloviev2024edaspy} and the
          \textsc{SMILE} inference engine~\cite{genie}.
          % , using the neonatal-jaundice ID of~\cite{gomez2007jaundice} as the benchmark.
    \item Empirically evaluate the impact of the main
          hyperparameters (population size, selection ratio, dead iterations, convergence criteria...) as well as different loss functions on the quality of the recovered parameters and
          on the run-time of the algorithm.
    \item Discuss the limits of the proposed approach and outline a set of natural extensions (LLM-informed initialisation,
          canonical-model parameterisation, parallel EDAs...).
\end{enumerate}

Through a simple model, \emph{``bypass''}, work is done to define the computational estimation methodology for the parameters, and it is tested with a complex real model, \emph{``NHLy''}.

% --------------------------------------------------------------
\section{Thesis Outline}
% --------------------------------------------------------------

The remainder of the document is organised in five further chapters.

\begin{itemize}
    \item \textbf{Chapter~\ref{ch:state-of-the-art}} reviews the state of the
          art: statistical learning from data, direct expert elicitation,
          evolutionary and generative approaches, and the inverse problem of
          recovering parameters from rules, positioning the contribution of this
          thesis.
    \item \textbf{Chapter~\ref{ch:framework}} develops the
          theoretical framework: Decision Support Systems, the
          formal definition of an Influence Diagram and its
          evaluation, the Knowledge-Acquisition problem that
          motivates this thesis, and a self-contained presentation
          of Estimation of Distribution Algorithms with the variants
          relevant for our setting.
    \item \textbf{Chapter~\ref{ch:implementation}} describes the
          implementation of the proof of concept: the architecture
          of the system, the encoding and decoding of the parameter
          vector, the design of the fitness function and the
          configuration of the EDA loop with \textsc{EDAspy}.
    \item \textbf{Chapter~\ref{ch:experiments}} reports the experimental study:
          the benchmark IDs, the protocol of the hyperparameter sweep,
          the comparison between the rules generated by the
          optimised model and those provided by the expert, and a
          discussion of the empirical results.
    \item \textbf{Chapter~\ref{ch:conclusions}} concludes the work, summarising the
          contributions and outlining the research directions
          identified in the state of the art that constitute
          natural lines of future work.
\end{itemize}
